% ==============================================================================
% PHẦN B: ĐI SÂU PHÂN TÍCH BÀI TOÁN LƯU TRỮ VẬT LÝ, CHI PHÍ CHỈ MỤC
% VÀ THỰC NGHIỆM ĐỐI CHỨNG TiDB VS. MySQL TRÊN DỮ LIỆU LỚN
% (Chiếm 50% nội dung báo cáo theo đúng định hướng của giảng viên)
% ==============================================================================

\phanbaocao{PHẦN B}{ĐI SÂU PHÂN TÍCH BÀI TOÁN LƯU TRỮ VẬT LÝ, CHI PHÍ CHỈ MỤC\\VÀ THỰC NGHIỆM ĐỐI CHỨNG TiDB VS. MySQL TRÊN DỮ LIỆU LỚN}

\section{Bài toán Tính toán Lý thuyết Dung lượng Lưu trữ Vật lý (Theoretical Storage Model)}
\label{sec:phan_b_tinh_toan_luu_tru}

\subsection{Đặc tả Lược đồ Bảng Thực nghiệm Chuẩn (\code{sbtest1})}

Để xây dựng mô hình toán học chuẩn xác đến từng byte và đối chiếu khách quan giữa kiến trúc phân trang B+Tree của MySQL InnoDB và cấu trúc Key-Value LSM-Tree của TiDB/TiKV, nghiên cứu sử dụng lược đồ bảng chuẩn công nghiệp \code{sbtest1} (từ bộ công cụ kiểm chuẩn Sysbench):

\begin{lstlisting}[language=SQL, caption={Cấu trúc bảng thực nghiệm chuẩn sbtest1}]
CREATE TABLE sbtest1 (
    id INT NOT NULL AUTO_INCREMENT,
    k INT NOT NULL DEFAULT '0',
    c CHAR(120) NOT NULL DEFAULT '',
    pad CHAR(60) NOT NULL DEFAULT '',
    PRIMARY KEY (id),
    KEY k_1 (k)
) ENGINE=InnoDB DEFAULT CHARSET=utf8mb4;
\end{lstlisting}

Bảng bao gồm 4 trường dữ liệu với các đặc tính:
\begin{itemize}
  \item \code{id}: Số nguyên 4 bytes có dấu, đóng vai trò Khóa chính (\code{PRIMARY KEY}).
  \item \code{k}: Số nguyên 4 bytes, có đánh Chỉ mục phụ (\code{KEY k\_1}).
  \item \code{c}: Chuỗi ký tự cố định 120 ký tự (chứa chuỗi ngẫu nhiên ASCII định dạng \code{utf8mb4}).
  \item \code{pad}: Chuỗi ký tự cố định 60 ký tự (chứa dữ liệu đệm ngẫu nhiên).
\end{itemize}

\subsection{Tính toán Kích thước Dữ liệu Thô Lý thuyết ($S_{\text{raw}}$)}

Kích thước dữ liệu thô thuần túy ($S_{\text{raw}}$) là dung lượng của các giá trị trường dữ liệu khi chưa tính đến bất kỳ siêu dữ liệu (\textit{metadata}), phần đệm trang hay cấu trúc chỉ mục bổ sung nào.

\begin{enumerate}
  \item \textbf{Kích thước 1 dòng dữ liệu thuần túy ($R_{\text{raw}}$):}
  \begin{equation}
    R_{\text{raw}} = \text{sizeof}(\code{id}) + \text{sizeof}(\code{k}) + \text{sizeof}(\code{c}) + \text{sizeof}(\code{pad})
  \end{equation}
  Thay số với dữ liệu ký tự ASCII chuẩn:
  \begin{equation}
    R_{\text{raw}} = 4\text{ B} + 4\text{ B} + 120\text{ B} + 60\text{ B} = \mathbf{188\text{ bytes/row}}
  \end{equation}

  \item \textbf{Kích thước thô của 1 mục chỉ mục phụ \code{KEY k\_1 (k)} ($I_{\text{raw}}$):}
  Chỉ mục phụ thô bao gồm giá trị khóa tìm kiếm \code{k} và con trỏ trỏ về bản ghi gốc (chính là khóa chính \code{id}):
  \begin{equation}
    I_{\text{raw}} = \text{sizeof}(\code{k}) + \text{sizeof}(\code{id}) = 4\text{ B} + 4\text{ B} = \mathbf{8\text{ bytes/entry}}
  \end{equation}

  \item \textbf{Quy mô 10.000.000 bản ghi ($10^7$ dòng):}
  \begin{itemize}
    \item Khi \textbf{CHƯA có Secondary Index \code{k\_1}}:
    \begin{equation}
      S_{\text{raw (no index)}} = 10^7 \times 188\text{ B} = 1.880.000.000\text{ B} \approx 1.792,9\text{ MiB} \approx \mathbf{1,75\text{ GiB}}
    \end{equation}
    \item Khi \textbf{ĐÃ CÓ Secondary Index \code{k\_1}}:
    \begin{equation}
      S_{\text{raw (with index)}} = 10^7 \times (188 + 8)\text{ B} = 1.960.000.000\text{ B} \approx 1.869,2\text{ MiB} \approx \mathbf{1,83\text{ GiB}}
    \end{equation}
  \end{itemize}

  \item \textbf{Mở rộng ở Quy mô 100.000.000 bản ghi ($10^8$ dòng --- Quy mô thử thách phần cứng 16GB RAM):}
  \begin{align}
    S_{\text{raw-100M (no index)}}   &= 10^8 \times 188\text{ B} \approx \mathbf{17,5\text{ GiB}}\\
    S_{\text{raw-100M (with index)}} &= 10^8 \times 196\text{ B} \approx \mathbf{18,3\text{ GiB}}
  \end{align}
  \textit{Nhận xét quan trọng:} Ở quy mô 100 triệu bản ghi, dung lượng dữ liệu thô ($18,3\text{ GiB}$) đã chính thức vượt quá tổng dung lượng RAM vật lý ($16\text{ GB}$) của máy chủ thử nghiệm. Điều này đặt hệ quản trị vào tình thế bắt buộc phải liên tục đọc/ghi và hoán đổi dữ liệu với ổ đĩa cứng.
\end{enumerate}

\subsection{Mô hình Lưu trữ Vật lý trên MySQL InnoDB (B+Tree 16KB)}

Trong MySQL InnoDB, dữ liệu bảng được tổ chức theo cấu trúc cây B+Tree cụm (\textit{Clustered Index}), trong đó các nút lá chứa trực tiếp toàn bộ dữ liệu dòng:

\begin{enumerate}
  \item \textbf{Cấu trúc vật lý của 1 dòng trong Clustered Index ($R_{\text{InnoDB}}$):}
  Mỗi bản ghi $188\text{B}$ khi ghi vào trang InnoDB phải gánh thêm các siêu dữ liệu quản lý:
  \begin{itemize}
    \item \textbf{Trường ẩn MVCC:} Gồm \code{DB\_TRX\_ID} ($6\text{ bytes}$) lưu mã giao dịch cuối cùng sửa đổi dòng và \code{DB\_ROLL\_PTR} ($7\text{ bytes}$) trỏ tới vùng Undo Log để phục vụ Rollback và đọc đa phiên bản $\rightarrow$ tốn thêm $13\text{ bytes}$.
    \item \textbf{Record Header:} Chứa cờ xóa mềm (\textit{deleted flag}), thứ tự bản ghi và con trỏ dòng tiếp theo $\rightarrow$ tốn $5\text{ bytes}$.
    \item \textbf{Column Length \& Null Offsets:} Mã hóa độ dài và giá trị rỗng $\rightarrow$ tốn $2\text{ bytes}$.
  \end{itemize}
  Tổng kích thước một bản ghi thực tế trong thân trang là:
  \begin{equation}
    R_{\text{InnoDB}} = 188\text{ B} + 13\text{ B} + 5\text{ B} + 2\text{ B} = \mathbf{208\text{ bytes/row}}
  \end{equation}
  Mỗi bản ghi cần thêm 1 con trỏ slot $2\text{ bytes}$ trong bảng thư mục chân trang (\textit{Page Directory}), nâng tổng chi phí lên $210\text{ bytes/row}$.

  \item \textbf{Độ phình do Phân mảnh Trang B+Tree 16KB (Page Overhead \& Fill Factor):}
  \begin{itemize}
    \item Mỗi trang dữ liệu chuẩn của InnoDB có kích thước $16\text{ KiB} = 16.384\text{ bytes}$, trong đó mất khoảng $120\text{ bytes}$ cho \code{File Header}, \code{Page Header} và \code{File Trailer}.
    \item Khi dữ liệu được nạp tuần tự tăng dần theo khóa chính \code{AUTO\_INCREMENT}, InnoDB chừa lại $1/16$ không gian trang trống dự phòng (hệ số lấp đầy tuần tự $F_{\text{seq}} = 15/16 = 93,75\%$).
    \item Số lượng dòng tối đa trên 1 trang lá được tính bằng:
    \begin{equation}
      N_{\text{rows/page}} = \left\lfloor \frac{16.384 \times \frac{15}{16} - 120}{210} \right\rfloor = \left\lfloor \frac{15.360 - 120}{210} \right\rfloor = \mathbf{72\text{ dòng/trang}}
    \end{equation}
    \item Số lượng trang lá cần thiết cho $10^7$ bản ghi:
    \begin{equation}
      N_{\text{leaf\_pages}} = \left\lceil \frac{10^7}{72} \right\rceil = \mathbf{138.889\text{ trang}}
    \end{equation}
    \item Tính thêm $1\%$ số trang nút nhánh nội bộ của cây B+Tree:
    \begin{equation}
      S_{\text{MySQL (chưa Index)}} \approx 138.889 \times 16\text{ KiB} \times 1,01 \approx 2.247.930\text{ KiB} \approx \mathbf{2.192\text{ MiB} \approx 2,14\text{ GiB}}
    \end{equation}
  \end{itemize}

  \item \textbf{Chi phí Lưu trữ khi Đánh thêm Secondary Index \code{k\_1} ($S_{\text{MySQL\_Index}}$):}
  \begin{itemize}
    \item Mỗi mục chỉ mục phụ trong cây B+Tree thứ cấp bao gồm: khóa \code{k} ($4\text{B}$) + con trỏ khóa chính \code{id} ($4\text{B}$) + Header chỉ mục ($6\text{B}$) + Slot Directory ($2\text{B}$) = $\mathbf{16\text{ bytes/entry}}$.
    \item Do giá trị cột \code{k} mang tính ngẫu nhiên, các trang B+Tree của chỉ mục phụ chịu hiện tượng tách trang (\textit{Page Split}) thường xuyên, dẫn đến hệ số lấp đầy thực tế chỉ đạt mức cân bằng B-Tree trung bình: $F_{\text{rand}} \approx \ln 2 \approx 70\% \sim 72\%$.
    \begin{equation}
      S_{\text{MySQL\_Index}} \approx \frac{10^7 \times 16\text{ B}}{0,72} \approx 222.222.222\text{ B} \approx \mathbf{212\text{ MiB} \approx 0,21\text{ GiB}}
    \end{equation}
  \end{itemize}

  \item \textbf{Tổng dung lượng vật lý file \code{sbtest1.ibd} của MySQL trên đĩa cứng:}
  \begin{align}
    S_{\text{MySQL (10M dòng có Index)}}  &= 2.192\text{ MiB} + 212\text{ MiB} = \mathbf{2.404\text{ MiB} \approx 2,35\text{ GiB}}\\
    S_{\text{MySQL (100M dòng có Index)}} &\approx \mathbf{24.040\text{ MiB} \approx 23,5\text{ GiB}}
  \end{align}
\end{enumerate}

\subsection{Mô hình Lưu trữ Vật lý trên Cụm TiDB / TiKV (RocksDB LSM-Tree \& Multi-Raft)}

Trong TiDB, dữ liệu không bị ràng buộc vào các trang 16KB cố định mà được nén thành các khối file SST trong cấu trúc LSM-Tree của RocksDB:

\begin{enumerate}
  \item \textbf{Kích thước Cặp Key-Value của Dòng Dữ liệu (Trước khi Nén):}
  \begin{itemize}
    \item \textbf{Khóa (\code{Key}):} Mã hóa \code{t{table_id}_r{row_id}} ($19\text{ bytes}$) + nhãn thời gian MVCC \code{commit\_ts} ($8\text{ bytes}$) + mã loại MemComparable ($1\text{ byte}$) = $\mathbf{28\text{ bytes}}$.
    \item \textbf{Giá trị (\code{Value}):} Dữ liệu 3 cột còn lại \code{k}, \code{c}, \code{pad} ($184\text{ bytes}$) + Header định dạng dòng TiDB Row Format v2 ($16\text{ bytes}$) + thông tin phiên bản MVCC trong Column Family \code{write} ($20\text{ bytes}$) = $\mathbf{220\text{ bytes}}$.
    \item Tổng kích thước một dòng chưa nén: $28 + 220 = \mathbf{248\text{ bytes/row}}$.
    \item $10^7$ dòng dữ liệu thô chưa nén: $10^7 \times 248\text{ B} \approx \mathbf{2.365\text{ MiB} \approx 2,31\text{ GiB}}$.
  \end{itemize}

  \item \textbf{Kích thước Cặp Key-Value của Secondary Index \code{k\_1} (Trước khi Nén):}
  \begin{itemize}
    \item \textbf{Khóa (\code{Key}):} \code{t{table_id}_i{index_id}_{k}_{id}} ($37\text{ bytes}$) + MVCC \code{commit\_ts} ($9\text{ bytes}$) = $\mathbf{46\text{ bytes}}$.
    \item \textbf{Giá trị (\code{Value}):} Trạng thái cờ trong CF \code{write} $\approx \mathbf{10\text{ bytes}}$.
    \item Tổng kích thước một mục chỉ mục phụ chưa nén: $46 + 10 = \mathbf{56\text{ bytes/entry}}$.
    \item $10^7$ mục chỉ mục chưa nén: $10^7 \times 56\text{ B} \approx \mathbf{534\text{ MiB} \approx 0,52\text{ GiB}}$.
    \item Tổng dung lượng chưa nén trên 1 bản sao (Dữ liệu + Chỉ mục): $2.365 + 534 = \mathbf{2.899\text{ MiB} \approx 2,83\text{ GiB}}$.
  \end{itemize}

  \item \textbf{Hiệu ứng Nén LSM-Tree ($C_{\text{ratio}}$) và Khuếch đại Không gian ($\alpha_{\text{LSM}}$):}
  Khác với B+Tree của MySQL để lãng phí không gian phân mảnh trang, RocksDB trong TiKV nén từng khối block $64\text{KB}$ trong SST bằng thuật toán \textbf{LZ4/Zstandard}.
  \begin{itemize}
    \item Nhờ cấu trúc chuỗi Key có chung tiền tố \code{t{table_id}} liên tục, tỷ lệ nén thực tế đo đạc được trên tập dữ liệu Sysbench đạt:
    \begin{equation}
      C_{\text{ratio}} \approx 0,60 \quad (\text{giảm } 40\% \text{ dung lượng thô})
    \end{equation}
    \item Hệ số khuếch đại không gian giữa các tầng $L_0 \rightarrow L_6$ trong mô hình Levelled Compaction của RocksDB là $\alpha_{\text{LSM}} \approx 1,11$ ($11\%$ dữ liệu phiên bản cũ đang chờ Compaction).
  \end{itemize}

  \item \textbf{Dung lượng thực tế trên 1 Node TiKV (1 Bản sao Dữ liệu):}
  \begin{itemize}
    \item Khi \textbf{CHƯA có Secondary Index}:
    \begin{equation}
      S_{\text{TiKV-1node (no index)}} = 2.365\text{ MiB} \times 0,60 \times 1,11 \approx \mathbf{1.575\text{ MiB} \approx 1,54\text{ GiB}}
    \end{equation}
    \textit{Nhận xét:} Trên 1 bản sao, TiDB chỉ tốn $1,54\text{ GiB}$, tiết kiệm hơn đáng kể so với con số $2,14\text{ GiB}$ của MySQL nhờ sức mạnh nén của LSM-Tree.
    \item Khi \textbf{ĐÃ CÓ Secondary Index \code{k\_1}}:
    \begin{equation}
      S_{\text{TiKV-1node (with index)}} = 2.899\text{ MiB} \times 0,60 \times 1,11 \approx \mathbf{1.930\text{ MiB} \approx 1,88\text{ GiB}}
    \end{equation}
  \end{itemize}

  \item \textbf{Tổng Dung lượng Toàn Cụm TiDB (3 Node TiKV --- Cơ chế Nhân bản 3 Replicas $N_{\text{rep}} = 3$):}
  Do TiDB đảm bảo tính sẵn sàng cao tuyệt đối bằng cách nhân bản mỗi Region thành 3 bản sao trên 3 máy chủ khác nhau, cộng thêm nhật ký đồng thuận \code{raftdb}:
  \begin{itemize}
    \item Khi \textbf{CHƯA có Secondary Index}:
    \begin{equation}
      S_{\text{TiDB-Cluster (no index)}} = 3 \times 1,54\text{ GiB} + 0,25\text{ GiB (Raft log)} \approx \mathbf{4,87\text{ GiB}}
    \end{equation}
    \item Khi \textbf{ĐÃ CÓ Secondary Index \code{k\_1}}:
    \begin{equation}
      S_{\text{TiDB-Cluster (with index)}} = 3 \times 1,88\text{ GiB} + 0,35\text{ GiB (Raft log)} \approx \mathbf{6,00\text{ GiB}}
    \end{equation}
    Dữ liệu được chia thành xấp xỉ $\mathbf{22\text{ Regions}}$ (mỗi Region kích cỡ trung bình $\sim 96\text{ MiB}$), rải đều sang 3 node TiKV (khoảng 7--8 Region Leader mỗi node).
    \item Với quy mô \textbf{100 triệu dòng}: Cụm TiDB chiếm khoảng $\mathbf{18,8\text{ GiB/node}}$ và tổng cộng $\mathbf{60\text{ GiB}}$ trên toàn cụm 3 máy chủ.
  \end{itemize}
\end{enumerate}

\subsection{Bảng Tổng hợp Đối chiếu Lý thuyết vs. Thực tế}

Bảng~\ref{tab:doi_chieu_dung_luong_tong_hop} tổng kết toàn bộ bức tranh dung lượng giữa mô hình dữ liệu thô, MySQL InnoDB và cụm phân tán TiDB.

\begin{table}[H]
\centering
\caption{Bảng đối chiếu tổng hợp dung lượng lưu trữ vật lý giữa MySQL InnoDB và TiDB}
\label{tab:doi_chieu_dung_luong_tong_hop}
\bangfont
\begin{tabularx}{\linewidth}{@{}>{\raggedright\arraybackslash}p{3.5cm}>{\centering\arraybackslash}p{2.3cm}>{\centering\arraybackslash}p{2.6cm}>{\centering\arraybackslash}p{2.7cm}C@{}}
\toprule
\textbf{Tiêu chí đối chiếu} & \textbf{Dữ liệu thô (Raw Data)} & \textbf{MySQL 8.0 (B+Tree 16KB)} & \textbf{TiDB (1 Node TiKV - RocksDB)} & \textbf{TiDB Toàn cụm (3 Node TiKV - Multi-Raft)}\\
\midrule
\textbf{Kích thước 1 dòng (Chưa Index)} & $188\text{ B}$ & $208\text{ B}$ (+20B MVCC/Header) & $248\text{ B}$ (thô) $\rightarrow \sim 165\text{ B}$ (sau nén) & $3 \times 165\text{ B} \approx 495\text{ B}$\\
\midrule
\textbf{Chi phí 1 mục Index \code{k\_1}} & $8\text{ B}$ & $16\text{ B}$ (chưa tính phân mảnh) & $56\text{ B}$ (thô) $\rightarrow \sim 37\text{ B}$ (sau nén) & $3 \times 37\text{ B} \approx 111\text{ B}$\\
\midrule
\textbf{10 Triệu dòng --- CHƯA Index} & \textbf{1,75 GiB} & \textbf{2,14 GiB} & \textbf{1,54 GiB} (Tiết kiệm hơn MySQL) & \textbf{4,87 GiB} (Đổi dung lượng lấy 3 bản sao HA)\\
\midrule
\textbf{10 Triệu dòng --- ĐÃ CÓ Index \code{k\_1}} & \textbf{1,83 GiB} & \textbf{2,35 GiB} (+212 MiB Index) & \textbf{1,88 GiB} (+355 MiB Index) & \textbf{6,00 GiB} (Chia thành $\sim 22$ Regions)\\
\midrule
\textbf{Giả định 100 Triệu dòng (Có Index)} & \textbf{18,3 GiB} & \textbf{23,5 GiB} (Vượt ngưỡng 16GB RAM) & \textbf{18,8 GiB/node} & \textbf{60,0 GiB} (Phân tán tải I/O trên 3 node)\\
\bottomrule
\end{tabularx}
\end{table}

% ==============================================================================
\section{Phân tích Chi phí Tính toán (RAM, CPU, IOPS) Trước và Sau khi có Index trên Máy 16GB RAM}
\label{sec:phan_b_chi_phi_tinh_toan}

\subsection{Kịch bản Đo đạc Thực nghiệm 1: Nạp 10 triệu bản ghi và Tạo Index}

Để phân tích chi phí phần cứng trên máy 16GB RAM, nhóm tiến hành đo đạc mức sử dụng tài nguyên thông qua script ghi log \code{docker stats} liên tục mỗi 2 giây qua 3 giai đoạn:

\subsubsection{Trạng thái 0: Mức chiếm dụng nền (Baseline trước khi nạp dữ liệu)}
\begin{itemize}
  \item Lượng RAM trống (\textit{Available RAM}) của hệ thống vật lý trước khi khởi động container đạt khoảng $13,8\text{ GB}$.
  \item Khi vừa khởi động:
  \begin{itemize}
    \item Container \code{mysql-single}: Chiếm khoảng $\mathbf{420\text{ MB}}$ RAM nền (do cấp phát cấu trúc khởi tạo của Buffer Pool).
    \item Cụm 5 containers TiDB: Tổng lượng RAM chiếm dụng ban đầu khoảng $\mathbf{1,25\text{ GB}}$ (trong đó mỗi node TiKV chiếm $\sim 300\text{ MB}$, TiDB Server $\sim 250\text{ MB}$, PD $\sim 100\text{ MB}$).
  \end{itemize}
\end{itemize}

\subsubsection{Trạng thái 1: Nạp 10 triệu bản ghi CHƯA có Secondary Index \code{k\_1}}
Tiến hành nạp dữ liệu tuần tự bằng lệnh \code{sysbench ... prepare}:
\begin{itemize}
  \item \textbf{Trên MySQL 8.0:}
  Do các dòng được chèn theo thứ tự \code{AUTO\_INCREMENT}, các trang lá B+Tree được ghi nối tiếp vào cuối cấu trúc cây. Quá trình này không phát sinh tách trang ngẫu nhiên. Mức RAM của \code{InnoDB Buffer Pool} tăng dần từ $420\text{ MB}$ lên chạm trần giới hạn cấu hình ($4\text{ GiB}$). Tốc độ nạp duy trì ở mức cao ($\sim 22.000\text{ rows/s}$), mức tiêu thụ CPU dao động $180\% \sim 240\%$ (tương đương 2--2,5 cores).
  \item \textbf{Trên TiDB:}
  Dữ liệu từ TiDB Server được nạp vào \code{MemTable} của TiKV. Khi dữ liệu vượt ngưỡng, quá trình tách vùng tự động diễn ra: từ $1\text{ Region}$ ban đầu, TiKV tự động băm nhỏ thành $\mathbf{18\text{ Regions}}$ dữ liệu và báo cáo về PD. PD điều phối rải đều các Region Leader và Follower sang 3 node TiKV. Bộ nhớ RAM của các node TiKV tăng lên (chủ yếu do nạp \code{BlockCache} chiếm $\sim 45\%$ RAM container). Tốc độ nạp đạt $\sim 16.500\text{ rows/s}$ (thấp hơn MySQL). Phép đo này không tách riêng phần của Raft, giao dịch hai pha, CPU và mạng.
\end{itemize}

\subsubsection{Trạng thái 2: Thực thi câu lệnh \code{ALTER TABLE sbtest1 ADD INDEX k\_1 (k)} trên 10 triệu dòng}
Đây là thời điểm bộc lộ rõ nét nhất sự khác biệt về chi phí tính toán giữa hai kiến trúc:
\begin{itemize}
  \item \textbf{Trên MySQL InnoDB (Kiến trúc B+Tree):}
  \begin{itemize}
    \item Quá trình \textit{Online DDL} của InnoDB buộc phải quét toàn bộ 10 triệu dòng trong Clustered Index, trích xuất cặp \code{(k, id)} và nạp vào vùng nhớ \code{innodb\_sort\_buffer\_size}.
    \item Do kích thước chỉ mục ($212\text{ MiB}$) vượt xa bộ nhớ sắp xếp, MySQL phải ghi các tập tin trung gian xuống đĩa tạm (\textit{Temporary Merge Sort on Disk}).
    \item Sau đó, InnoDB tiến hành xây dựng cây B+Tree từ dưới lên (\textit{Bulk Index Build}). Mức tải I/O đĩa tăng vọt lên mức tối đa của ổ cứng ($120 \sim 180\text{ MB/s}$ ghi đĩa liên tục), CPU bị khóa ở mức $100\%$ của 1 core xử lý phân loại.
    \item \textbf{Nguy cơ khi chèn dữ liệu vào bảng đã có Index:} Nếu bảng đã có sẵn \code{KEY k\_1}, mỗi thao tác chèn bản ghi ngẫu nhiên buộc InnoDB phải tìm kiếm và nạp trang lá tương ứng của \code{k\_1} từ đĩa lên Buffer Pool. Khi bảng phình lên 100 triệu dòng, Buffer Pool bị tràn, dẫn đến hiện tượng tráo trang liên tục (\textit{Buffer Pool Thrashing}) và bão I/O đọc/ghi ngẫu nhiên.
  \end{itemize}

  \item \textbf{Trên TiDB (Kiến trúc Distributed Reorg \& RocksDB LSM-Tree):}
  \begin{itemize}
    \item TiDB sử dụng thuật toán thêm chỉ mục bất đồng bộ phân tán (\textbf{Distributed Online DDL}). Node TiDB khởi tạo các luồng Reorg Workers (\code{tidb\_ddl\_reorg\_worker\_cnt = 4}), chia nhỏ 18 Regions thành các đoạn khóa con (\code{batches}) và giao cho 3 node TiKV đọc song song.
    \item Các bản ghi chỉ mục \code{t\_i\_k\_id} được đóng gói thành các cặp Key-Value và đẩy ngược vào TiKV dưới dạng các lệnh ghi Key-Value thông thường.
    \item \textbf{Đặc điểm ghi của LSM-Tree:} Thay vì phải tìm trang B+Tree cũ trên đĩa để sửa đổi (\textit{Read-Modify-Write}), RocksDB trên TiKV chỉ đơn giản là nạp cặp Key-Value mới vào \code{MemTable} trên RAM và ghi tuần tự vào nhật ký \code{WAL}. Đường ghi vào \code{WAL} là tuần tự, nhưng các lần flush và \textit{compaction} về sau ghi lại dữ liệu nhiều lần, tạo ra khuếch đại ghi và tải đĩa nền.
  \end{itemize}
\end{itemize}

\subsection{So sánh Chi phí Truy vấn và Kế hoạch Thực thi (\code{EXPLAIN ANALYZE})}

Nhóm khảo sát câu lệnh tìm kiếm theo giá trị cột chỉ mục:
\begin{lstlisting}[language=SQL]
EXPLAIN ANALYZE SELECT * FROM sbtest1 WHERE k = 500000;
\end{lstlisting}

Bảng~\ref{tab:so_sanh_explain} đối chiếu chi phí thực thi trước và sau khi đánh chỉ mục trên cả hai hệ thống.

\begin{table}[H]
\centering
\caption{Đối chiếu kế hoạch thực thi và tài nguyên tiêu tốn trước và sau khi có Index \code{k\_1}}
\label{tab:so_sanh_explain}
\bangfont
\begin{tabularx}{\linewidth}{@{}>{\raggedright\arraybackslash}p{2.8cm}>{\raggedright\arraybackslash}p{4.8cm}Y@{}}
\toprule
\textbf{Trường hợp} & \textbf{MySQL 8.0 InnoDB} & \textbf{TiDB Distributed NewSQL}\\
\midrule
\textbf{CHƯA có Index} & \textbf{Toán tử:} \code{Table Full Scan}\par
- Đọc toàn bộ $138.889\text{ trang}$ ($2,14\text{ GiB}$).\par
- Quét đủ $10.000.000\text{ dòng}$ qua CPU.\par
- Làm dội sạch toàn bộ Buffer Pool.\par
- Thời gian thực thi: $\mathbf{\sim 3,82\text{ giây}}$. & \textbf{Toán tử:} \code{TableFullScan} đẩy xuống \code{TiKV Coprocessor}\par
- 3 node TiKV quét song song các Region trên đĩa.\par
- Trả về dòng thỏa mãn lên TiDB Server.\par
- \code{total\_process\_keys}: $10.000.000\text{ keys}$.\par
- Thời gian thực thi: $\mathbf{\sim 1,45\text{ giây}}$ (nhanh hơn nhờ quét song song 3 node).\\
\midrule
\textbf{ĐÃ CÓ Index \code{k\_1}} & \textbf{Toán tử:} \code{Index Lookup} trên B+Tree \code{k\_1}\par
- Đọc 3--4 tầng trang B+Tree.\par
- Lấy con trỏ \code{id} và truy cập Clustered Index.\par
- Đọc đúng 1 trang đĩa.\par
- Thời gian thực thi: $\mathbf{\sim 0,38\text{ ms}}$. & \textbf{Toán tử:} \code{IndexLookUp} (2 chặng phân tán)\par
- Chặng 1: \code{IndexRangeScan} quét đúng khóa chỉ mục trong \code{CF write} trên TiKV để lấy \code{row\_id}.\par
- Chặng 2: \code{TableRowIDScan} lấy dữ liệu đầy đủ.\par
- \code{total\_process\_keys}: $1 \sim 2\text{ keys}$.\par
- Thời gian thực thi: $\mathbf{\sim 1,12\text{ ms}}$ (thời gian toàn bộ câu lệnh, chưa tách thành phần).\\
\bottomrule
\end{tabularx}
\end{table}

\textit{Kết luận chuyên môn:}
Khi có Index, MySQL đạt độ trễ thấp hơn TiDB trong phép thử này. Thời gian tổng của câu lệnh \textbf{không đủ để quy chênh lệch cho mạng}. Khoảng chênh gồm nhiều thành phần: lượt quay về bảng của \code{IndexLookUp} (hiệu ứng của index không covering), xử lý và xếp hàng của cop task, đọc MVCC và LSM-Tree, mã hóa gRPC và truyền dữ liệu. Mục~\ref{sec:phuong_phap_tach_chi_phi} trình bày cách nhóm tách các thành phần này. Khi chưa có Index, TiDB đẩy phép quét xuống Coprocessor trên 3 node TiKV.

% ==============================================================================
\section{Thiết kế thực nghiệm chi phí chỉ mục phụ trên dữ liệu vận đơn}
\label{sec:phan_b_thuc_nghiem_doi_chung}

\subsection{Thực nghiệm sơ bộ và động lực nghiên cứu}

Ở giai đoạn 1, nhóm chạy sysbench trên \code{sbtest1} 10 triệu dòng với 16 đến 512 luồng. MySQL đạt thông lượng cao hơn TiDB 13--32 lần ở tải \code{oltp\_read\_write} và 8--13 lần ở tải \code{oltp\_update\_index}. Giả thuyết ban đầu rằng TiDB ổn định vượt trội ở 256--512 luồng \textbf{không được số liệu xác nhận}. Khoảng chênh này bao gồm TSO, giao dịch hai pha, Raft, giới hạn CPU của từng TiKV và bộ nhớ đệm. Phần còn lại của nghiên cứu hỏi: bao nhiêu trong đó gắn với việc duy trì và truy cập chỉ mục phụ?

\subsection{Phương pháp tách hiệu ứng index và chi phí truy cập KV}
\label{sec:phuong_phap_tach_chi_phi}

Nhóm không gán mọi chênh lệch độ trễ cho mạng. Mỗi hiệu ứng được đo bằng một phép so sánh riêng, trong cùng một hệ, trước khi so hai hệ.

\begin{table}[H]
\centering
\caption{Các hiệu ứng cần tách và phép so sánh tương ứng}
\label{tab:tach_hieu_ung}
\bangfont
\begin{tabularx}{\linewidth}{@{}>{\raggedright\arraybackslash}p{3.6cm}YY@{}}
\toprule
\textbf{Hiệu ứng} & \textbf{Phép so sánh} & \textbf{Bằng chứng}\\
\midrule
Covering & Cùng câu SQL, cùng tiền tố khóa, có và không có cột bao phủ & Toán tử truy cập, p50/p95/p99, số key\\
Tính duy nhất (unique) & \code{idx\_code} (unique) so với \code{idx\_code\_nonuniq} & \code{Point\_Get} so với quét index\\
Độ chọn lọc & RQ2: ép index, ép quét bảng, optimizer tự chọn & Điểm giao nhau, điểm đổi kế hoạch\\
Truy cập KV phân tán & Số cop task, số RPC, \code{rpc\_time}, thời gian xử lý và chờ & \code{EXPLAIN ANALYZE}\\
Nhân bản Raft & RF=3 so với RF=1, cùng 3 TiKV & TPS, p95/p99, byte ghi mỗi giao dịch\\
\bottomrule
\end{tabularx}
\end{table}

Mọi container nằm trên một máy. Do đó \code{rpc\_time} gồm cả xử lý phía TiKV, xếp hàng và gRPC, không phải độ trễ mạng thuần. So sánh RF=3 với RF=1 đo tổng chi phí nhân bản (Raft log, fsync, CPU của Follower và tranh chấp trên ổ đĩa chung), không phải riêng phần mạng.

\subsection{RQ1: covering index trên truy vấn vận đơn}

\begin{itemize}
  \item \textbf{Q1:} tra một mã vận đơn. So \code{idx\_code}, \code{idx\_code\_nonuniq} và \code{idx\_code\_cov} trên cùng câu \code{SELECT status, updated\_at}. Cặp hai index không unique cho biết hiệu ứng covering, cặp unique và không unique cho biết hiệu ứng \code{Point\_Get}.
  \item \textbf{Q2:} 20 đơn mới nhất của một shop trong 7 ngày. So \code{idx\_sender}, \code{idx\_sender\_created} và \code{idx\_sender\_cov}. Đây là bằng chứng chính cho covering vì ba index dùng cùng câu SQL và đều không unique.
  \item \textbf{Q3:} đếm đơn theo trạng thái tại một kho. So \code{idx\_hub} và \code{idx\_hub\_status}.
\end{itemize}

Với mỗi trường hợp, nhóm ghi kế hoạch thực thi, p50/p95/p99, số dòng trả về, số cop task, số RPC, \code{rpc\_time} và dung lượng index. Chi phí ghi của covering được đo ở RQ4.

\subsection{RQ4: chi phí ghi}

Ma trận chính gồm 4 bộ index (0, 1, 3, 5), 2 mức luồng (16, 64), 2 tải (W1 chèn đơn, W2 cập nhật \code{status} và \code{updated\_at}), 3 lần lặp, trên MySQL, TiDB RF=3 và TiDB RF=1. Thêm 24 lượt so \code{idx\_sender\_created} với \code{idx\_sender\_cov} ở 16 luồng trên MySQL và TiDB RF=3. W2 chỉ làm thay đổi \code{idx\_sender\_cov}, vì cột \code{status} nằm trong index này.

Kết quả W2 phải diễn giải theo \textbf{số index bị ảnh hưởng}: \code{idx\_code}, \code{idx\_sender\_created} và \code{idx\_cod} không chứa cột được cập nhật, còn \code{idx\_hub\_status} và \code{idx\_status} thì có. Dung lượng index của TiDB trong \code{information\_schema} là ước lượng từ thống kê, không dùng làm dung lượng vật lý.

% TODO(TV2, TV3): điền số liệu RQ1, RQ2, RQ4 sau khi có CSV chính thức.

\subsection{Demo bổ sung: chịu lỗi và mở rộng}

Phần này chỉ là demo bổ sung từ giai đoạn 1 trên \code{sbtest1}, không phải kết quả chính. Khi \code{tikv0} bị \code{docker kill} ở giây thứ 60 dưới tải \code{oltp\_read\_write} 64 luồng, QPS của TiDB bằng 0 trong 10 giây và về lại 90\% sau 18,4 giây mà không cần can thiệp. 0/2.035 lệnh ghi probe bị lỗi. Với MySQL một node, QPS bằng 0 trong 46 giây, cần khởi động lại và về lại 90\% sau 48,3 giây, 140/2.050 lệnh probe bị lỗi. Cả hai không mất giao dịch đã commit. Phép so sánh không ngang hàng vì MySQL không có replica. Nó minh họa khác biệt về tính sẵn sàng, không phải về độ bền.

Khi thêm một TiKV lúc đang có tải, Region đầu tiên chuyển sang node mới sau 16 giây và ứng dụng không gặp lỗi. TPS giảm 25\% vì node mới không mang thêm phần cứng thật. Demo này cho thấy khả năng mở rộng trực tuyến, không chứng minh thêm node làm tăng hiệu năng.
